\documentclass[11pt,a4paper]{article}
\usepackage{amsmath,amssymb,amsthm,mathtools}
\usepackage{geometry}
\usepackage{hyperref}
\usepackage{booktabs}
\usepackage{enumitem}
\usepackage{xcolor}
\usepackage{seqsplit}

\geometry{margin=1in}
\hypersetup{colorlinks=true,linkcolor=blue,citecolor=blue,urlcolor=blue}

\newtheorem{theorem}{Theorem}[section]
\newtheorem{lemma}[theorem]{Lemma}
\newtheorem{proposition}[theorem]{Proposition}
\newtheorem{corollary}[theorem]{Corollary}
\newtheorem{definition}[theorem]{Definition}
\newtheorem{remark}[theorem]{Remark}

\title{\textbf{Metriplectic Completion of the Frohmanian Symplectic Tether:\\
An Assembled Finite-Dimensional Viscous System,\\ the Imported Tether Kernel, and Lean~4 Checks}}
% Author: Benjamin Stanley Frohman
\author{Benjamin Stanley Frohman\\
Independent Researcher, Texas, United States\\
ORCID: 0009-0006-7068-3718\\
\texttt{frohmanbenjamin@gmail.com}}
\date{Version 1.0, November 2026}

\sloppy
\emergencystretch=3em
\begin{document}
\maketitle

\begin{abstract}
We assemble a concrete metriplectic system with viscous dissipation and test the Frohmanian tether
kernel against the metriplectic axioms, with every result checked by the Lean~4 kernel. The system
lives on \(\mathbb{R}^n\times\mathbb{R}\) (velocity coefficients \(u\) and an entropy coordinate \(s\)): a Poisson
bracket from a state-dependent skew matrix \(J(u)\), a viscous metric bracket with viscosities
\(\nu_i\ge0\), energy \(H=\tfrac12\|u\|^2+s\) and entropy \(S=s\). We prove all metriplectic axioms, derive the
equations of motion \(\dot u_i=(J(u)u)_i-\nu_iu_i\), \(\dot s=\sum_i\nu_iu_i^2\), and show non-vacuity for a
three-mode triad instance (nonzero brackets, nonconstant energy, strictly positive entropy production,
nonzero reversible motion). The tether kernel is imported from its own Lean development at a pinned
commit, not re-proved. Against it we prove: the metric bracket of any metriplectic system with the
tether's kinetic-energy Hamiltonian vanishes against \(H\) in both slots; a tether summand next to a
degenerate viscous summand is forced to satisfy the (unproved) degeneracy condition (C2); and the tether
kernel is negative semidefinite and, by the imported nondegeneracy theorem, strictly negative
somewhere, so \emph{no} metriplectic system has the tether kernel as its metric bracket. The assembled
system is finite-dimensional and is not derived from a Fourier truncation of the Navier--Stokes
equations; nothing here recovers or formalizes the continuum Navier--Stokes or Navier--Stokes--Fourier
equations. This is not a solution of the Clay Millennium problem, and a kernel check is not a
geometric or physical proof.
\end{abstract}

\tableofcontents

\section{Introduction}
The Euler equations of ideal incompressible hydrodynamics are Lie--Poisson dynamics on the dual of
the Lie algebra of divergence-free vector fields (Arnold, 1966). Adding viscosity gives the
Navier--Stokes equations, which are dissipative. The metriplectic formalism
\cite{Morrison1986,MorrisonUpdike2024} adds a symmetric positive-semidefinite \emph{metric bracket}
to a Poisson bracket so that energy is conserved and entropy does not decrease.

The Frohmanian tether kernel \(T\) \cite{Frohman2026} is a symmetric correction to the Arnold
bracket. Its Lean development records an obstruction: if \(T\) were degenerate on the kinetic-energy
Hamiltonian (condition (C2)), symmetry would force \(T(H,M)=0\) for every \(M\), so the tether
would supply no damping along the energy direction. Condition (C2) is \emph{not} proved for \(T\)
in that development.

This paper asks what the metriplectic formalism says about the tether, and records exactly what has
been checked by the Lean kernel. Section~\ref{sec:lean} lists the Lean declarations and how they were
audited. The scope is algebraic and finite-dimensional; continuum statements in
Sections~\ref{sec:arnold}--\ref{sec:metriplectic} are background, cited from the literature, and are not
formalized.

\section{Arnold Lie--Poisson structure (background)}\label{sec:arnold}
Let \(\mathfrak{g} = \mathfrak{X}_{\mathrm{div}}(\mathbb{T}^3)\) be the Lie algebra of smooth
divergence-free vector fields on the three-torus, with dual identified with vorticity fields. The
Lie--Poisson bracket of two functionals of the vorticity \(\omega\) is
\begin{equation}
\{F,G\}(\omega) = \int_{\mathbb{T}^3} \omega \cdot \Bigl(
\frac{\delta F}{\delta\omega}\times\frac{\delta G}{\delta\omega}\Bigr)\,dx,
\end{equation}
and the kinetic energy \(H(\omega) = \frac12\int_{\mathbb{T}^3}|u|^2\,dx\), \(u = \mathrm{curl}^{-1}\omega\),
generates the Euler equations \(\partial_t\omega + \mathcal{L}_u\omega = 0\). This section is
classical background and is not formalized here.

\section{Metriplectic systems}\label{sec:metriplectic}
\begin{definition}[Metriplectic system]\label{def:ms}
A metriplectic system on a state set \(X\) consists of a Poisson bracket \(\{\cdot,\cdot\}\) and a metric
bracket \((\cdot,\cdot)\) on functionals \(X\to\mathbb{R}\), an energy \(H\) and an entropy \(S\), such that
for all functionals \(F,G\) and all states:
\begin{align}
\{F,G\} &= -\{G,F\} && \text{(skew-symmetry)},\label{eq:skew}\\
(F,G) &= (G,F) && \text{(symmetry)},\label{eq:sym}\\
(F,F) &\ge 0 && \text{(positive semidefiniteness)},\label{eq:pos}\\
\{F,S\} &= 0 && \text{(entropy is a Casimir)},\label{eq:cas}\\
(F,H) &= 0 && \text{(energy degeneracy)}.\label{eq:null}
\end{align}
The evolution of an observable is \(\dot F = \{F,H\} + (F,S)\).
\end{definition}
The Jacobi identity for \(\{\cdot,\cdot\}\) is not part of this definition and is not used below.
The previous version of this paper listed ``\((S,H)\ge 0\)'' as the entropy condition; that was an
error, since \((S,H)=0\) by \eqref{eq:null} and \eqref{eq:sym}. The correct condition is
\eqref{eq:pos}, which gives \((S,S)\ge0\).

For the compressible Navier--Stokes--Fourier system, Morrison, Updike, Zaidni and Barham construct
the metric bracket from a 4-bracket built from viscosity and conductivity operators
\cite{MorrisonUpdike2024,ZaidniMorrisonBenjelloun2024,BarhamMorrisonZaidni2024}. We cite those
constructions as background only. \textbf{They are not formalized in this work, and no statement
here recovers the continuum equations.}

\section{Thermodynamic consistency (formalized)}\label{sec:thermo}
\begin{theorem}[Energy conservation]\label{thm:energy}
In a metriplectic system, \(\dot H = \{H,H\} + (H,S) = 0\).
\end{theorem}
\begin{proof}
By \eqref{eq:skew} with \(F=G=H\), \(\{H,H\}=-\{H,H\}\), so \(\{H,H\}=0\). By \eqref{eq:sym},
\((H,S)=(S,H)\), and \((S,H)=0\) by \eqref{eq:null} with \(F=S\).
\end{proof}

\begin{theorem}[Entropy production]\label{thm:entropy}
In a metriplectic system, \(\dot S = \{S,H\} + (S,S) = (S,S) \ge 0\).
\end{theorem}
\begin{proof}
By \eqref{eq:skew}, \(\{S,H\} = -\{H,S\}\). By \eqref{eq:cas} applied to the functional \(F=H\),
\(\{H,S\}=0\). Hence \(\{S,H\}=0\) and \(\dot S=(S,S)\), which is \(\ge0\) by \eqref{eq:pos} with \(F=S\).
\end{proof}

\begin{remark}
The draft Lean proof of Theorem~\ref{thm:entropy} cited the Casimir condition at \(F=S\), which gives
\(\{S,S\}=0\), not the needed \(\{S,H\}=0\), and did not compile. The Lean proof now follows the argument
above exactly: Casimir condition at \(F=H\), then skew-symmetry.
\end{remark}

\begin{theorem}[Symmetric swap]\label{thm:swap}
If \(B\) is symmetric and \(B(F,H)=0\) for every \(F\), then \(B(H,F)=0\) for every \(F\).
\end{theorem}
\begin{proof}
\(B(H,F)=B(F,H)=0\).
\end{proof}

\section{The Frohmanian tether and the damping obstruction}\label{sec:tether}
In the imported Lean development \cite{Frohman2026}, the tether kernel is defined as
\begin{equation}
T(F,G)(\omega) = -\kappa\,\|\omega\|_{L^2}^2\int_{\mathbb{T}^3}
\bigl(\Pi_u\tfrac{\delta F}{\delta\omega}\bigr)\cdot
\bigl(\Pi_u\tfrac{\delta G}{\delta\omega}\bigr)\,dx,\qquad \|\omega\|_{L^2}^2=\int|\omega|^2,
\end{equation}
with the global enstrophy outside the integral, \(\kappa>0\), and \(\Pi_u\) the \(L^2\)-orthogonal
projector onto the complement of \(\mathrm{span}\{u\}\), \(u\) the velocity recovered from \(\omega\).
The previous version of this paper wrote a pointwise weight \(|\omega(x)|^2\) inside the integral;
that is a different object and is not the one formalized. The three conditions are:
\begin{enumerate}[label=(C\arabic*)]
\item coadjoint invariance --- \textbf{not proved} for \(T\) in that development;
\item degeneracy on the kinetic-energy Hamiltonian, \(T(F,H)=0\) for all \(F\) --- \textbf{not proved}
for \(T\) in that development (an earlier route assumed \(\delta H/\delta\omega=u\), which is false with
its definitions, and was removed);
\item controllable negative feedback, \(T(F,F)=-\kappa\|\omega\|^2\int\|\Pi_u\delta F\|^2\) --- proved, as an
equality.
\end{enumerate}

\paragraph{Correction to the previous version.} The previous version said that ``a theorem already
verified in Lean'' gives \(T(M,H)=0\) for every \(M\). That was a misattribution: \(T(M,H)=0\) for all
\(M\) \emph{is} condition (C2), a hypothesis. What the Lean development proves is the other slot,
\(T(H,M)=0\), \emph{assuming} (C2) (lemma \texttt{left\_H\_zero\_of\_C2\_symm}, specialized in
\texttt{tetherKernel\_left\_H\_zero}). Since (C2) is not proved for \(T\), neither identity is
established unconditionally for \(T\).

\section{The obstruction stated against the imported kernel}\label{sec:bridge}
The Lean file \texttt{MetriplecticNS/TetherBridge.lean} imports the tether development and works
on its state space (\texttt{CoadjointOrbit}) with its Hamiltonian (\texttt{KineticEnergyHamiltonian}).

\begin{theorem}[Metric obstruction]\label{thm:bridge}
Let \((\{\cdot,\cdot\},(\cdot,\cdot),H,S)\) be a metriplectic system on divergence-free vorticity fields with
\(H\) equal to the tether's kinetic-energy Hamiltonian. Then the metric bracket satisfies (C2), and
\((H,M)=0\) for every functional \(M\).
\end{theorem}
\begin{proof}
(C2) for the metric bracket is \eqref{eq:null}. The imported lemma
\texttt{left\_H\_zero\_of\_C2\_symm}, applied to the metric bracket with symmetry \eqref{eq:sym} and this
(C2), gives \((H,M)=0\).
\end{proof}

\begin{theorem}[Tether summand]\label{thm:summand}
In the setting of Theorem~\ref{thm:bridge}, suppose \((F,G) = V(F,G) + T(F,G)\) for some bracket \(V\) with
\(V(F,H)=0\) for all \(F\). Then \(T\) satisfies (C2), and \(T(H,M)=T(M,H)=0\) for every \(M\).
\end{theorem}
\begin{proof}
For every \(F\): \(0=(F,H)=V(F,H)+T(F,H)=T(F,H)\), which is (C2) for \(T\). The imported theorem
\texttt{tetherKernel\_left\_H\_zero} then gives both identities.
\end{proof}
So, under the hypotheses of Theorem~\ref{thm:summand}, the tether summand does not couple to the energy
direction at all. Theorem~\ref{thm:summand} does not prove (C2) for \(T\); it shows that (C2) for \(T\) is
\emph{forced} by placing \(T\) in such a split metric bracket.

\begin{lemma}[Sign of the tether diagonal]\label{lem:diag}
\(T(F,F)(\omega)\le0\) for every \(F\) and \(\omega\).
\end{lemma}
\begin{proof}
By (C3), \(T(F,F)(\omega)=(-\kappa)\cdot\|\omega\|^2\cdot\int\|\Pi_u\delta F\|^2\). Here \(-\kappa\le0\) since \(\kappa>0\);
\(\|\omega\|^2=\int|\omega|^2\ge0\); and the last integral has a nonnegative integrand, so it is \(\ge0\) (in
Lean, the Bochner integral of a nonnegative function is nonnegative, including when the function is not
integrable and the integral is \(0\) by convention). A nonpositive number times two nonnegative numbers is
nonpositive.
\end{proof}

\begin{theorem}[Sign obstruction]\label{thm:sign}
If the metric bracket of a metriplectic system on vorticity fields equals \(T\), then \(T(F,F)=0\) for all
\(F\) and all \(\omega\), and in particular \(\dot S=(S,S)=0\): no entropy is produced.
\end{theorem}
\begin{proof}
\eqref{eq:pos} gives \(T(F,F)\ge0\); Lemma~\ref{lem:diag} gives \(T(F,F)\le0\). Then \(\dot S=(S,S)\) by
Theorem~\ref{thm:entropy}, and \((S,S)=T(S,S)=0\).
\end{proof}
\begin{corollary}[Non-vacuity: the tether is not a metric bracket]\label{cor:nometric}
There are \(\omega,F\) with \(T(F,F)(\omega)<0\). Consequently no metriplectic system on vorticity fields
has metric bracket equal to \(T\).
\end{corollary}
\begin{proof}
The imported theorem \texttt{tetherKernel\_nondegenerate} gives a regular vorticity \(\omega\) (a smooth
compactly supported rotation field) and a linear test functional \(F\) with \(T(F,F)(\omega)\neq0\);
Lemma~\ref{lem:diag} gives \(T(F,F)(\omega)\le0\), hence \(<0\). If a metriplectic system had metric bracket
\(T\), Theorem~\ref{thm:sign} would give \(T(F,F)(\omega)=0\), a contradiction.
\end{proof}
So Theorem~\ref{thm:sign} is not vacuous: its hypothesis is never satisfied, and that is a proved fact
about the actual tether kernel. Thus a positive entropy production cannot come from the tether kernel. In a split \((\cdot,\cdot)=V+T\),
\((S,S)=V(S,S)+T(S,S)\le V(S,S)\): the tether can only lower the entropy production supplied by \(V\), and
\(V\) must dominate it for \eqref{eq:pos} to hold. Whether any metriplectic system of the split form of Theorem~\ref{thm:summand} exists with \(T\neq0\) is
open: by that theorem, existence would prove (C2) for \(T\). These statements replace the previous version's claim
that ``the tether kernel may be inserted as one summand of the metric bracket'' and that ``energy
conservation continues to hold because \(T\) is already degenerate on \(H\)''; that degeneracy is not proved.

\section{A finite-dimensional viscous metriplectic system}\label{sec:galerkin}
This section is the concrete instance formalized in \texttt{MetriplecticNS/Galerkin.lean}. It is a
genuine metriplectic system, but it is finite-dimensional, and it is not derived from the Navier--Stokes
equations.

\subsection{Definition}
Fix \(n\in\mathbb{N}\). The state is \(x=(u,s)\in\mathbb{R}^n\times\mathbb{R}\) (\(u\) are velocity coefficients, \(s\) is an
entropy / internal-energy coordinate). For a functional \(F:\mathbb{R}^n\times\mathbb{R}\to\mathbb{R}\) let
\(\partial_{u_i}F(x)=DF(x)[(e_i,0)]\) and \(\partial_sF(x)=DF(x)[(0,1)]\), where \(DF(x)\) is the Fréchet
derivative (taken to be \(0\) where \(F\) is not differentiable, Mathlib's convention, so that the brackets
are defined on every functional). The data are a map \(J:\mathbb{R}^n\to\mathbb{R}^{n\times n}\) with
\(J(u)^{\mathsf T}=-J(u)\) for every \(u\), and viscosities \(\nu_i\ge0\). Define
\begin{align*}
\{F,G\}(x)&=\sum_{i,j}\partial_{u_i}F\;J(u)_{ij}\;\partial_{u_j}G, &
(PF)_i(x)&=\partial_{u_i}F-(\partial_sF)\,u_i,\\
(F,G)(x)&=\sum_i\nu_i\,(PF)_i\,(PG)_i, &
H(u,s)&=\tfrac12\|u\|^2+s,\qquad S(u,s)=s.
\end{align*}
(For a Fourier truncation one would take \(\nu_i=\nu|k_i|^2\); we do not need that interpretation.)

\subsection{Derivatives}
\begin{lemma}\label{lem:derivs}
For every state \(x=(u,s)\) and index \(i\):
\(\partial_{u_i}H=u_i\), \(\partial_sH=1\); \(\partial_{u_i}S=0\), \(\partial_sS=1\); for \(K=\tfrac12\|u\|^2\),
\(\partial_{u_i}K=u_i\), \(\partial_sK=0\); for the coordinate \(c_j(x)=u_j\), \(\partial_{u_i}c_j=\delta_{ij}\),
\(\partial_sc_j=0\). Consequently \((PH)_i=0\), \((PS)_i=-u_i\), \((PK)_i=u_i\), \((Pc_j)_i=\delta_{ij}\).
\end{lemma}
\begin{proof}
\(x\mapsto\|u\|^2\) has derivative \(v\mapsto2\langle u,v_1\rangle\) (chain rule: \(\|\cdot\|^2\) has derivative
\(2\langle u,\cdot\rangle\) and \(x\mapsto u\) is the continuous linear projection). Hence
\(DK(x)[v]=\langle u,v_1\rangle\) and \(DH(x)[v]=\langle u,v_1\rangle+v_2\), since \(x\mapsto s\) is the continuous linear
projection onto the second factor. Evaluating at \(v=(e_i,0)\) gives \(\langle u,e_i\rangle=u_i\); at \(v=(0,1)\)
gives \(0\) for \(K\) and \(1\) for \(H\). \(S\) and \(c_j\) are continuous linear maps, so they are their own
derivatives: \(DS(x)[v]=v_2\) and \(Dc_j(x)[v]=(v_1)_j\). The values of \(P\) follow from the definition, e.g.
\((PH)_i=u_i-1\cdot u_i=0\).
\end{proof}

\subsection{The structure axioms}
\begin{theorem}\label{thm:axioms}
The data above satisfy \eqref{eq:skew}--\eqref{eq:null}.
\end{theorem}
\begin{proof}
\emph{Skew-symmetry.} Exchange the order of summation and use \(J_{ji}=-J_{ij}\):
\(\{G,F\}=\sum_{i,j}\partial_{u_i}G\,J_{ij}\,\partial_{u_j}F=\sum_{i,j}\partial_{u_j}G\,J_{ji}\,\partial_{u_i}F
=-\sum_{i,j}\partial_{u_i}F\,J_{ij}\,\partial_{u_j}G=-\{F,G\}\).
\emph{Symmetry} of \((\cdot,\cdot)\) is commutativity of the product \((PF)_i(PG)_i\).
\emph{Positivity:} \((F,F)=\sum_i\nu_i(PF)_i^2\ge0\) since \(\nu_i\ge0\).
\emph{Casimir:} \(\partial_{u_j}S=0\) for every \(j\), so every term of \(\{F,S\}\) vanishes.
\emph{Energy degeneracy:} \((PH)_i=0\) for every \(i\), so every term of \((F,H)\) vanishes.
\end{proof}

\subsection{Equations of motion}
\begin{proposition}\label{prop:eom}
For \(\dot F=\{F,H\}+(F,S)\):
\[
\dot u_i=(J(u)u)_i-\nu_iu_i,\qquad \dot s=\sum_i\nu_iu_i^2,\qquad
\frac{d}{dt}\tfrac12\|u\|^2=-\sum_i\nu_iu_i^2,\qquad \dot H=0,\qquad \dot S\ge0.
\]
\end{proposition}
\begin{proof}
For \(F=c_i\): \(\{c_i,H\}=\sum_{a,b}\delta_{ia}J_{ab}u_b=(Ju)_i\) and
\((c_i,S)=\sum_a\nu_a\delta_{ia}(-u_a)=-\nu_iu_i\).
For \(F=S\): \(\{S,H\}=0\) since \(\partial_{u}S=0\), and \((S,S)=\sum_a\nu_a(-u_a)^2\).
For \(F=K\): \(\{K,H\}=\sum_{a,b}u_aJ_{ab}u_b=0\), because the quadratic form \(Q=u^{\mathsf T}Ju\) of a skew
matrix satisfies \(Q=-Q\) (exchange the summation order and use skew-symmetry); and
\((K,S)=\sum_a\nu_au_a(-u_a)\). The last two claims are Theorems~\ref{thm:energy} and~\ref{thm:entropy}
applied to the system of Theorem~\ref{thm:axioms}.
\end{proof}
So kinetic energy lost to viscosity equals entropy gained: \(\frac{d}{dt}K=-\dot S\).

\subsection{A three-mode triad instance and non-vacuity}
Take \(n=3\), \(\nu_i=1\), and
\[
J(u)=\begin{pmatrix}0&u_2&u_1\\-u_2&0&3u_0\\-u_1&-3u_0&0\end{pmatrix},\qquad
J(u)u=(2u_1u_2,\;2u_0u_2,\;-4u_0u_1),
\]
a quadratic interaction with \(2+2-4=0\), of the energy-conserving triad type that appears in Galerkin
truncations of fluid equations. Skew-symmetry of \(J(u)\) is checked entry by entry. We do not claim the
Jacobi identity for this \(J\), and we do not claim this triad is a Fourier truncation of Navier--Stokes.

\begin{proposition}[Non-vacuity]\label{prop:nonvac}
Let \(x_0=(e_1+e_2,0)\). Then: \(\{c_0,c_1\}(x_0)=1\neq0\); \((c_0,c_0)(x_0)=1\neq0\); \(H(0,0)=0\neq1=H(0,1)\);
\(\dot S(x_0)=2>0\); and \(\dot u_0(x_0)=2u_1u_2-u_0=2\), so the reversible part of the motion is nonzero.
\end{proposition}
\begin{proof}
In general \(\{c_i,c_j\}=J(u)_{ij}\) and \((c_i,c_i)=\nu_i\) by Lemma~\ref{lem:derivs}. At \(u=e_1+e_2=(0,1,1)\):
\(J_{01}=u_2=1\); \(\nu_0=1\); \(\dot S=\sum u_i^2=2\); \(\dot u_0=2\cdot1\cdot1-0=2\).
\end{proof}

\section{Lean~4 formalization and reproducibility}\label{sec:lean}
\paragraph{Modules.}
\begingroup\raggedright
\begin{itemize}
\item \texttt{MetriplecticNS/Basic.lean}: \texttt{MetriplecticSystem}; \texttt{energy\_conserved},
\texttt{entropy\_production}, \texttt{entropy\_production\_eq}, \texttt{tether\_damping\_obstruction}
(Theorems~\ref{thm:energy}--\ref{thm:swap}).
\item \texttt{MetriplecticNS/TetherBridge.lean}: imports \texttt{FrohmanianTether.Combined};
\texttt{metric\_degenerate}, \texttt{metric\_left\_H\_zero}, \texttt{tetherKernel\_summand\_no\_H\_coupling},
\texttt{tetherKernel\_diag\_nonpos}, \texttt{tetherKernel\_as\_metric\_diag\_zero},
\texttt{tetherKernel\_as\_metric\_no\_entropy\_production}, \texttt{tetherKernel\_diag\_neg\_somewhere},
\texttt{no\_metriplectic\_system\_with\_tether\_metric} (Section~\ref{sec:bridge}); it also imports
\texttt{FrohmanianTether.TetherNondegeneracy}.
\item \texttt{MetriplecticNS/Galerkin.lean}: the system \texttt{system D}, the five axioms, equations of
motion \texttt{evolution\_coord}, \texttt{evolution\_S}, \texttt{evolution\_K},
\texttt{kinetic\_loss\_eq\_entropy\_gain}, and the triad instance with \texttt{triad\_poisson\_ne\_zero},
\texttt{triad\_metric\_ne\_zero}, \texttt{H\_not\_const}, \texttt{triad\_entropy\_production\_pos},
\texttt{triad\_evolution\_coord0} (Section~\ref{sec:galerkin}).
\end{itemize}
\endgroup

\paragraph{Dependency.} The tether development is the Lake package \texttt{FrohmanianTether} from
\url{https://github.com/BenFrohman/frohmanian-tether-sdiff}, pinned to commit
\texttt{\seqsplit{8b7bd6741f8a54587db6a95b4a5199dcf2775f97}}. Both repositories use Lean
\texttt{v4.30.0-rc1} and Mathlib commit \texttt{\seqsplit{0c154d67103f74be3a0f2c509f72ccbf5be9f2a7}}, so the kernel
lemmas are imported rather than re-proved.

\paragraph{Audit.} \texttt{scripts/kernel\_check.sh} rejects \texttt{sorry}, \texttt{admit},
\texttt{native\_decide}, \texttt{axiom}, \texttt{opaque} and \texttt{unsafe} in the sources, and runs
\texttt{\#print axioms} on all 75 declarations of the three modules. Each depends only on
\texttt{propext}, \texttt{Classical.choice} and \texttt{Quot.sound} (this output is transitive, so it covers
the imported tether lemmas). Continuous integration runs \texttt{lake build} and the audit on every push
and uploads the audit output.

\paragraph{Errors caught by Lean checking.} Kernel checking caught real errors in the original
upload: the draft entropy-production proof used the wrong rule (\(\{S,S\}=0\), the Casimir condition at
\(F=S\)) where the argument needs \(\{S,H\}=0\) (Casimir condition at \(F=H\) plus skew-symmetry), and the
Navier--Stokes ``specialization'' was a placeholder (state type \texttt{Unit}, both brackets identically
zero) that was never assembled into a system. Both are corrected here: the proof uses the right rule,
and the placeholder is replaced by the assembled system of Section~\ref{sec:galerkin}.

\paragraph{Rebuild.} \texttt{lake exe cache get \&\& lake build \&\& ./scripts/kernel\_check.sh}.
The exact source commit for this version is \texttt{\seqsplit{1ecea6cf99011cd4b962e68ea05a1a03a90e3eb3}}.

\paragraph{What is not formalized.} The Arnold bracket on vorticity fields, functional derivatives on
Sobolev spaces, any integration by parts, the viscous term \(\nu\Delta u\), the 4-bracket constructions, and
recovery of the Navier--Stokes or Navier--Stokes--Fourier equations. The previous version said the viscous
term was recovered ``by integration by parts'' in Lean and that the specialization to integral
functionals of fluid densities was formalized; neither was true, and neither is claimed now.

\section{Conclusion and open problems}
The metriplectic formalism sharpens the earlier obstruction. With the tether's own Hamiltonian, any
metriplectic metric bracket vanishes against \(H\) in both slots; a tether summand next to a degenerate
viscous summand is forced to satisfy (C2) and then vanishes against \(H\); and the tether's diagonal has the
wrong sign to produce entropy, so positive entropy production must come from another summand, such as a
viscous one; and since the tether is strictly negative somewhere, it is not the metric bracket of any
metriplectic system. A concrete finite-dimensional viscous metriplectic system, assembled and checked
in Lean, shows the structure is consistent and non-vacuous.

Open problems:
\begin{itemize}
\item prove or refute (C2) for the tether kernel;
\item derive a finite-dimensional system of the type of Section~\ref{sec:galerkin} from an actual Fourier
truncation of the Navier--Stokes nonlinearity, with a proof that the truncated bracket is skew;
\item formalize an infinite-dimensional metriplectic structure for the incompressible equations on Sobolev
spaces;
\item find a non-symmetric or higher-order geometric correction that evades the obstruction while keeping
coadjoint equivariance.
\end{itemize}
Nothing here bears directly on the global regularity of the three-dimensional Navier--Stokes equations.

\begin{thebibliography}{9}
\bibitem{Arnold1966}
V.~I. Arnold.
\newblock Sur la géométrie différentielle des groupes de Lie de dimension infinie et ses applications à l'hydrodynamique des fluides parfaits.
\newblock \emph{Ann. Inst. Fourier} \textbf{16} (1966), 319--361.

\bibitem{Morrison1986}
P.~J. Morrison.
\newblock A paradigm for joined Hamiltonian and dissipative systems.
\newblock \emph{Physica D} \textbf{18} (1986), 410--419.

\bibitem{Materassi2014}
M. Materassi.
\newblock Metriplectic algebra for dissipative fluids in Lagrangian formulation.
\newblock arXiv:1409.6542, 2014.

\bibitem{MorrisonUpdike2024}
P.~J. Morrison and M.~H. Updike.
\newblock Inclusive curvature-like framework for describing dissipation: metriplectic 4-bracket dynamics.
\newblock \emph{Phys. Rev. E} \textbf{109} (2024), 045202.

\bibitem{ZaidniMorrisonBenjelloun2024}
A. Zaidni, P.~J. Morrison and S. Benjelloun.
\newblock Thermodynamically consistent Cahn--Hilliard--Navier--Stokes equations using the metriplectic dynamics formalism.
\newblock arXiv:2402.11116, 2024.

\bibitem{BarhamMorrisonZaidni2024}
W. Barham, P.~J. Morrison and A. Zaidni.
\newblock A thermodynamically consistent discretization of 1D thermal-fluid models using their metriplectic 4-bracket structure.
\newblock arXiv:2410.11045, 2024.

\bibitem{Frohman2026}
B.~S. Frohman.
\newblock The Frohmanian Tether on \(\mathrm{SDiff}(\mathbb{T}^3)\): Lean development
(\texttt{frohmanian-tether-sdiff}, commit \texttt{8b7bd67}).
\newblock Zenodo, 2026. DOI: 10.5281/zenodo.24120781.
\end{thebibliography}

\end{document}
